\section{Introdução}
\label{sec:intro}

Trinta e cinco milhões de brasileiros --- 35\% da força de trabalho ocupada --- atuam em ocupações com exposição quase nula à inteligência artificial ($\theta \le 2$ em uma escala de 0 a 10). Não é a sofisticação das tarefas que as isola, mas o caráter predominantemente manual, presencial e informal de suas atividades. O trabalho doméstico (6{,}7 milhões de postos) e a construção civil estrutural (3{,}2 milhões) somam mais empregos do que todas as ocupações de tecnologia, finanças e administração somadas; sua exposição ao índice de \citet{eloundou2024gpts} é de 0{,}1 e 0{,}9, respectivamente, contra 7{,}5 dos desenvolvedores de software e 8{,}6 dos operadores de máquinas de escritório.

Esse quadro qualifica a narrativa dominante sobre IA e mercado de trabalho. Nos Estados Unidos, \citet{eloundou2024gpts} estimam que cerca de 80\% da força de trabalho tem ao menos 10\% de suas tarefas afetadas por grandes modelos de linguagem. No Brasil, a mesma métrica revela um padrão assimétrico: a exposição é \emph{positivamente} associada à escolaridade e à renda, concentrando-se no emprego formal. O mercado de trabalho brasileiro conta com um amortecedor ausente nas economias avançadas --- a informalidade ---, e esse amortecedor situa-se exatamente na base do espectro de exposição.

Este artigo investiga quem, no mercado de trabalho brasileiro, está exposto à IA e por meio de quais mecanismos econômicos. Desenvolvo um modelo de designação de tarefas com setor informal em que a IA individual está universalmente disponível, mas a IA corporativa exige capital organizacional restrito à formalidade. Testo as previsões do modelo com os microdados trimestrais da PNAD Contínua (2026Q1, com mais de 227 mil trabalhadores analisados individualmente), combinados ao mapeamento determinístico entre a COD brasileira e a classificação O*NET-SOC.

Três resultados empíricos emergem. Primeiro, a exposição ordena-se fortemente com a escolaridade: com controles mincerianos completos (idade, idade ao quadrado, gênero e raça) e erros-padrão agrupados por ocupação (122 grupos), cada ano adicional de estudo associa-se a 0{,}23 ponto a mais de exposição à IA (erro-padrão 0{,}03; $p < 0{,}001$; $N = 227.629$). Segundo, a associação entre exposição e formalidade é marcante: na relação incondicional, cada ponto de exposição associa-se a $-6{,}23$ p.p. de informalidade ($p < 0{,}001$). Condicional à escolaridade individual, o coeficiente da exposição atenua-se em $37{,}2\%$ para $-3{,}91$ p.p. (erro-padrão 0{,}96), enquanto cada ano de escolaridade associa-se a $-2{,}61$ p.p. de informalidade ($p < 0{,}001$), evidenciando a mediação parcial pelo capital humano e a persistência de um canal tecnológico direto. Terceiro, o gradiente educacional é mais acentuado nas Unidades da Federação com maior maturidade do setor formal: a interação entre escolaridade e a formalidade da UF (calculada via \emph{leave-one-out}) é de 0{,}28 (erro-padrão 0{,}06; $p < 0{,}001$).

A interpretação desses padrões é de \emph{sorting} ocupacional de equilíbrio, não de impacto causal ou deslocamento imediato de trabalhadores. O modelo formaliza essa distinção: trabalhadores mais escolarizados possuem vantagem comparativa em tarefas cognitivas e abstratas complementares à IA, auto-selecionando-se nas ocupações mais expostas do setor formal, onde o retorno ao capital organizacional corporativo compensa os custos de conformidade tributária e regulatória.

O artigo conecta três literaturas. A primeira mensura a exposição das ocupações a novas tecnologias \citep{frey2017future,acemoglu2022tasks,eloundou2024gpts}, desenvolvida primordialmente para economias formais; dentro dela, a automação desloca e reinstata tarefas \citep{acemoglu2019automation}, e o emprego persiste porque máquinas complementam --- e não apenas substituem --- trabalho \citep{autor2015why}. A segunda analisa a escolha entre formalidade e informalidade em economias em desenvolvimento \citep{ulyssea2018firms}. A terceira modela a alocação de trabalhadores a tarefas \citep{acemoglu2011skills,autordorn2013growth}.

A contribuição do trabalho é tripla: (i) apresenta a primeira mensuração em nível de microdados da exposição à IA para o Brasil com controles demográficos completos e inferência robusta; (ii) formaliza a barreira do capital organizacional corporativo como mediador da exposição tecnológica sob informalidade; e (iii) qualifica o debate de políticas públicas, indicando que as políticas de adaptação e capacitação em IA devem focar no estrato médio formal escolarizado, hoje o mais exposto aos choques tecnológicos.

A Seção~\ref{sec:modelo} apresenta o modelo teórico; a Seção~\ref{sec:dados} descreve os microdados e o cruzamento ocupacional; a Seção~\ref{sec:estrategia} detalha a estratégia empírica; a Seção~\ref{sec:resultados} discute as estimativas e testes de robustez; e a Seção~\ref{sec:conclusao} conclui.
